\subsection{The partial stationary trace}\label{c4:partial-trace}

When the two color blocks have different sizes, their suppression
profiles turn off at different distances from the central mark.  Between
those distances only the larger block is suppressed.  We must estimate
the stationary trace in this region before estimating the full product
of radial links. In the stationary trace on either sector of slopes (character \(\sigma=0,2v\) in \eqref{c4:PR-a}--\eqref{c4:PR-d}), suppress \(1F,12,E3\) of color sector 1, leaving all other screening activities unchanged. The resulting normalized trace \(I_\sigma(L)\), including \eqref{c4:PR-d}, satisfies
\begin{equation}
                    |I_\sigma(L)|\ \le\ \exp\big((5/288+o(1))L\big)
                    \qquad(L\to+\infty).                                 \label{c4:PS-1}
\end{equation}
The proof reduces the surviving two-root expansion to the positive
root lattice \(D_4\).  A finite Lie-algebra calculation then determines a
unitary character of a rank-two invariant lattice.  Poisson summation on
that lattice supplies the exponential rate.

\paragraph{Reducing the partially suppressed trace.} Suppression has exactly the same meaning (omit the corresponding exponential factors) before the radial transformation. Using \eqref{c4:Pair-Bcut} and removing \(g_2\) cyclically gives angular twist
\[
 \exp(-\mathrm{i}c E_{E4}^{(1)})\exp(c E_{24}^{(1)})
       \exp(2\pi i(\sigma,P))D_1\widehat S_1\widehat C.
\]
Indeed the two surviving factors shown on the left commute with the conjugation, and the indicated character can be moved through all the root actions used here. Put \(B'=S_\alpha C\), \(\alpha=\alpha_{10}\), to avoid using \(B\) for both twist and vertex. These manipulations can first be made power by power in the angular energy and evaluated for small positive \(Q_L\), by the cut finiteness and spread estimates (the \(B\) cut of that estimate is \(B'\)); equivalently expand the two root exponentials at the new cut, bunching their contours closely there. At a fixed cut energy each product has finite transitions on any input, identities of factors hold as operators by the angular calculation, and in cyclic cancellation of \(g_2\) the diagonal support is finite by the stated bound for the \(B'\) cut including \(g_2^{\pm1}\). On the radial side deleting factors still gives real-analytic convergence along \(0<L<\infty\).

Write \(Z_T=\prod_{k\ge1}\det(1-Q_L^k T)^{-1}\) with determinant on the pertinent charge space. In Euclidean four-space take
\[
\begin{gathered}
 D=\{x\in\mathbb Z^4:\ \sum x_i\in 2\mathbb Z\},\\
 (f_1,f_2,f_3,f_4)=(e_1+e_2,e_1+e_3,e_1+e_4,e_1-e_3),\\
 B_4=\tfrac12\begin{pmatrix}1&1&-1&1\\-1&1&1&1\\-1&1&-1&-1\\1&1&1&-1\end{pmatrix}
\end{gathered}
\]
in orthonormal coordinates. The \(f\)'s form an integral basis of \(D\); \(B_4\) preserves \(D\), with characteristic polynomial \(\Phi_{12}\) (write \(\Phi_j\) for cyclotomic polynomials); \(K_4=(B_4-I)^{-1}\) preserves \(D\) and its dual since \(\det(B_4-I)=1\). Use the ordinary Fock states and fields on this lattice with cocycle \(\exp(\pi i(b,K_4P))\), always \(b\in D\), and \(X_b\) the zero coefficient for roots \(b^2=2\). Lift \(B_4\) by pure charge and oscillator transformation without a scalar. With \(U=\exp X_{f_1}\exp X_{f_2}\widehat B_4\), we claim
\begin{equation}
 I_\sigma(L)=
 \frac{Z_{B'}}{Z_C\sum_{k\in\mathbb Z}Q_L^{(4kp_C)^2/2}}\,
              \frac{\operatorname{Tr}_{D}(UQ_L^H)}{Z_{B_4}} .           \label{c4:PS-2}
\end{equation}
Here the trace on \(D\) uses \textbf{lattice} momenta (we can still use the same operators on all dual momenta).

We verify both the contractions and the scalar phases in
\eqref{c4:PS-2}; matching the oscillator determinant alone would not
identify the trace. The two charges \(u,z\), in the order of factors from left to right, are \(\alpha_{17}-\alpha_{11}-\alpha_{12},-\alpha_{12}\); in computations below they map to \(f_1,f_2\), intertwining the orbit pairings for \(B'\) with those for \(B_4\). Indeed for \(T=(u,z)\),
\((B')^6 T=-T\), and the \(2\times2\) Gram blocks against \((B')^j T,\ j=0,\ldots,5\), in row-major tuples, are
\[
 (2,1,1,2),\ (1,0,1,-1),\ (1,1,0,1),\ (0,-1,1,0),\
 (-1,0,-1,-1),\ (-1,-1,0,1),
\]
exactly the Euclidean ones. Further \(B'\) has characteristic polynomial
\(\Phi_4^2\Phi_8\Phi_6\Phi_{12}\) (use that of \(C\), \((x^8-1)^2/(x^2-1)\), times \(1+(\alpha,C(x-C)^{-1}\alpha)\)). Thus in each expanded term with counts \(r,s\) of \(u,z\), the pre-twist momentum \(m\) is uniquely \((1-B')^{-1}(r u+s z)\); it always lies in \(\mathcal L^*\). The Euclidean calculation likewise uses \(m'=(1-B_4)^{-1}(r f_1+s f_2)\) in \(D\). All resulting momentum powers (including heat) and oscillator pair/self contractions agree by the orbit pairings, even if the first orbit span has a radical. The free oscillator determinants are those indicated in \eqref{c4:PS-2}.

For extra detail on the remaining scalars, let \(M=(1-B')^{-1}T\). Multiplication using the angular Gram and \(J\) gives its two columns (first-sector coefficients, second zero)
\[
 (1,0,-1,-1,-1,-1,-1,0),\qquad (1,1,-1,-1,-1,-1,-1,-1)/2
\]
with integral pairings as claimed. With \(H_4=h_1+C^2\alpha-\alpha\), the other entries needed are
\[
 \begin{gathered}
 (\alpha,CM)=(-1,-1),\quad J(\alpha,h_1)=-1,\quad J(\alpha,CM)=(3,2),
 \quad J(T,B'M)/2=\begin{pmatrix}9/2&2\\2&1\end{pmatrix},\\
 (v,B'M)=0,\qquad J(u,z)=1,\qquad J(T,H_4)=(0,1)^t .
 \end{gathered}
\]
Thus the total coefficient phase in units \(\pi i\) is (put \(b=r u+s z,\ k=(\alpha,Cm)\))
\[
\begin{aligned}
 &\tfrac12[kJ(\alpha,h_1-Cm)+k^2-k]-(\alpha,B'm)/4\\
 &\qquad{}+J(b,B'm)/2+rs/2-J(b,H_4)/2-5r/4-3s/4\\
 &\quad=13r^2/2+8rs+5s^2/2-(r+s)/2 .
\end{aligned}
\]
On the Euclidean side \((f_i,K_4 f_j)_{i,j\le2}=\left(\begin{smallmatrix}-1&0\\-1&-1\end{smallmatrix}\right)\), and \((r f_1+s f_2,K_4 B_4 m')-\binom r2-\binom s2\) gives \(11r^2/2+6rs+3s^2/2+(r+s)/2\). Their difference is \((r+s)^2-(r+s)\), an even integer. This proves the reduction using the angular trace expansion at small \(Q_L\) (or formal zero modes). The closed Euclidean expression below is real-analytic for all real \(L>0\); thus \eqref{c4:PS-2} extends there by the radial convergence.

\paragraph{A finite Lie-algebra calculation.} We now calculate the Euclidean trace. Here \(D\) is the root lattice of type \(D_4\), but we include details of the automorphism argument. Its 24 roots are the vectors \(\pm e_i\pm e_j,\ i<j\), equivalently the list
\[
 \mathcal R=(B_4^j f_l:0\le j<6,\ l=1,2)\ \text{ followed by its negatives}
\]
(\(l\) runs fastest). The Cartan operators \((h,P)\) together with \(X_b,\ b\in\mathcal R\), close into a Lie algebra of dimension at most 28: the residue calculation of the angular method applies verbatim with \(K_4\) since \(K_4+K_4^*=-I\) and the pairings are integral. In particular \([X_b,X_{-b}]=-(b,P)\); for \(b+d\) a root, \([X_b,X_d]=(-1)^{(b,K_4 d)}X_{b+d}\), with zero for other nonopposites. Commutation by each zero-mode generator has the same action as this adjoint calculation on the corresponding \textbf{full fields}, replacing \((h,P)\) by its Heisenberg current and \(X_b\) by \(W_b\); this follows from the same residue and oscillator brackets.

On momenta \(e_1+D\) at \(H=1/2\) there are just the eight states \(\pm e_i\); Cartans are diagonal and \(X_b\) takes \(m\to m+b\) with coefficient \((-1)^{(b,K_4m)}\) if possible. These matrices preserve infinitesimally the nondegenerate symmetric form pairing \(m,-m\) with coefficient 1 (use \((b,K_4 b)=-1\)). They are independent by their supports, so identify the Lie algebra exactly with the corresponding \(\mathfrak{so}(8,\mathbb C)\). Exponentials of any Lie elements are defined also on the whole Fock algebra, in finite energy spaces since now the form is positive.

The next finite calculation supplies the Cartan subalgebra and the two
trace values from which we will determine the invariant character.
\begin{lemma}\label{c4:partial-automorphism}
The adjoint action of \(U\) on \(\mathfrak{so}(8,\mathbb C)\) fixes a
regular semisimple element \(x\) whose eigenvalues in the
eight-dimensional vector representation are purely imaginary. Moreover,
\[
 \operatorname{tr}\operatorname{Ad}(U)=1,\qquad
 \operatorname{tr}\operatorname{Ad}(U^2)=-1.
\]
\end{lemma}
\begin{proof}
Only the Lie element \(x\), not \(U\), is represented on the eight
weight states. In order Cartans \(f_i\) then \(\mathcal R\), take
\[
 x=(-1,1,0,1; -1,-2,0,-2,0,-1,0,-1,0,-1,0,-1; 0,-1,0,-1,0,-1,0,-1,0,-2,1,-3).
\]
Write \(A_0\) acting by \(B_4\) on Cartans and the unweighted \(B_4\)-permutation on roots. Write \(J_0=\operatorname{ad}(X_{f_1}+X_{f_2})\) by the brackets just given; \(J_0^3=0\), equivalently by direct application (the two commuting roots also belong to a root triple system of type \(A_2\), where their sum of operators is conjugate to one root). For \(T_i=J_0^i A_0/i!,\ i=0,1,2\), multiplication gives
\[
 \operatorname{tr}T_i=(0,1,0)_i,\qquad
 (\operatorname{tr}T_iT_j)=\begin{pmatrix}2&-2&0\\-2&1&0\\0&0&0\end{pmatrix}.
\]
In the same basis \(T_1x=(0,1,0,0;\ -1,2,1,0,0,1,0,0,0,0,0,0; 0,0,0,0,0,0,0,0,0,-1,1,-1)\),
\(T_2x\) has just first two root entries \(-1,-1\), and \(A_0 x=x-T_1x-T_2x\). The sum of the \(T_i\) is the required action. Finally the matrix of \(x\) (positive weights \(e_i\) then negative) and its determinant polynomial are
\[
 \begin{pmatrix}
1&0&0&-1&0&-1&2&0\\
0&-1&-3&1&1&0&0&1\\
0&1&0&1&-2&0&0&-2\\
1&-1&-2&0&0&-1&2&0\\
0&0&1&0&-1&0&0&-1\\
0&0&0&0&0&1&-1&1\\
-1&0&0&-1&0&3&0&2\\
0&0&1&0&1&-1&-1&0
 \end{pmatrix},\qquad
 \lambda^2(\lambda^2+8)(\lambda^4+12\lambda^2+4)
\]
(for instance the even traces of powers \(2,4,6,8\) are \(-40,400,-4192,45120\)). The six nonzero eigenvalues are distinct purely imaginary and the remaining kernel is two-dimensional (skew operators for the form have even rank). Thus \(x\) is semisimple. Its Cartan coordinates consist of three nonzero values, distinct up to sign, and one zero value; every root pairing, a sum or difference of two coordinates, is nonzero. This proves regularity as well as the stated trace identities.
\end{proof}

\paragraph{The invariant lattice and its character.}
The fixed regular element allows us to read the infinite Fock trace from
the action on a Cartan subalgebra. All conjugations and cyclic trace identities in this step are first performed on each finite-dimensional \(H\)-eigenspace of the positive \(D_4\) lattice. The zero modes preserve these spaces, so no boundedness of a conjugating operator on the completed Fock space is needed. We sum over energies only after identifying the unitary character below, when the resulting oscillator and theta series converge absolutely.

An orthogonal conjugation of determinant one puts \(x\) in the Cartan as an element \(y\) with nonzero purely imaginary pairings against all roots: choose paired eigenvectors on opposite eigenspaces, a paired isotropic basis on the kernel, adjusting orientation by interchanging a pair if necessary. This conjugation is a product of exponentials in the Lie representation and hence can be performed also on the Fock operators using zero modes.

To recall the elementary argument, an orthogonal matrix for a nondegenerate complex symmetric form is a product of reflections in nonisotropic vectors, by matching an orthonormal basis successively (for two unit vectors, their sum or their difference is nonisotropic, allowing at most two reflections). Thus an even product is joined to identity by deforming each reflecting line to a common one, possible along complex paths avoiding zero norm. Near identity the logarithm is skew by orthogonality, giving the claimed product. Bracket compatibility ensures that lifting any such Lie exponentials gives the required action on the currents at once.

After conjugation the new twist fixes \(y\), hence preserves the Cartan (its centralizer) and permutes root spaces, inducing an orthogonal root permutation \(W\) also preserving the positive system cut out by the imaginary pairings. This same action holds on the current fields. Thus the twist sends an empty oscillator state \(m\) to \(t_m|Wm\rangle\) by energy and Cartan weights, with \(t_m\ne0,\ t_0=1\) (the zero modes and lattice lift used here preserve the vacuum), and acts compatibly by \(W\) on the oscillators. It permutes the classes of dual momenta modulo \(D\) exactly as \(B_4\) since root rotations preserve those classes. In particular the permutation has order divisible by 3: \(B_4 e_1,B_4^2 e_1,e_1\) are in distinct classes and \(B_4^3 e_1\) is in \(e_1+D\). On the other hand \(W\) acts by a permutation of the three outer simple roots of the positive system, fixing the central one. Indeed signed permutation of coordinates puts the defining regular vector for positives in the usual chamber \(a_1>a_2>a_3>|a_4|\); the positive roots there have as indecomposables \(e_1-e_2,e_2-e_3,e_3-e_4,e_3+e_4\), with the second central. Thus \(W\) cycles the three outer ones.

Write \(\beta\) for that central simple root, \(\gamma\) for the sum of the three outer ones. The invariant sublattice is \(\mathbb Z\beta+\mathbb Z\gamma\), with \(\gamma^2=6,\ \beta^2=2,\ (\gamma,\beta)=-3\); its roots are \(\pm\beta,\pm(\gamma+\beta),\pm(\gamma+2\beta)\). For each invariant root \(b\), if its adjoint eigenvalue is \(c_b\), full field covariance implies \(t_{m+b}=c_b t_m\) on invariant \(m\), simply by the pure-output coefficient of \(W_b\). Hence \(t\) there is a character, generated using \(b=\beta,\gamma+\beta\), and \(c_b=t_b\). Put \(\lambda_i=t_b\) for the three positive invariant roots in displayed order; \(\lambda_3=\lambda_1\lambda_2\). Only fixed roots contribute to the Lie traces of the first two powers (Cartan trace 1 both times), so \(p_i=\lambda_i+\lambda_i^{-1}\) have \(\sum p_i=0,\ \sum p_i^2=4\) by the checked traces. Since \(\sum p_i^2-p_1p_2p_3=4\), their values are \(0,\pm\sqrt2\). In particular the character is unitary and the absolute argument residues of \(\lambda_1,\lambda_2,\lambda_3\) modulo 8 in units \(2\pi/8\) are \(1,2,3\) in some order.

We have proved
\[
 \operatorname{Tr}_D(UQ_L^H)
     =Z_W\sum_{m\in\mathbb Z\gamma+\mathbb Z\beta} t_m Q_L^{m^2/2}.
\]
\paragraph{The exponential rate.}
Put \(\Lambda=\mathbb Z\beta+\mathbb Z(\gamma+\beta)\). In this basis
its Gram matrix is \(\left(\begin{smallmatrix}2&-1\\-1&2\end{smallmatrix}\right)\),
so its covolume is \(\sqrt3\). Since the character is unitary, choose a
real vector \(\vartheta\) with
\(t_m=e^{2\pi i(\vartheta,m)}\) on \(\Lambda\). Ordinary Gaussian
Poisson summation, using \(Q_L=e^{-4\pi^2/L}\), gives
\[
 \sum_{m\in\Lambda}t_m Q_L^{m^2/2}
   =\frac{L}{2\pi\sqrt3}
        \sum_{q\in\Lambda^*}e^{-L|q-\vartheta|^2/2}.
\]
The right side has positive terms, so its exponential rate is determined
by the nearest point of this dual coset, without cancellation.
Its minimum squared norm is the minimum of
\((B^2+BC+C^2)/96\), where \(B,C\) range over representatives of the
angle residues of \(\lambda_1,\lambda_2\); this is \(7/96\).
For nearest representatives the absolute residues of the sum enforce
equal sign for \(1,2\), opposite sign when one is 3; other lifts have a
coordinate at least 5 in modulus, and
\(B^2+BC+C^2\ge 3\max(B^2,C^2)/4>7\) for those lifts.
Thus the momentum growth rate in \(\log(\cdot)/L\) is \(-7/192\);
the Poisson prefactor is only polynomial.

Finally the rates for \(Z_{B'},Z_C,Z_W,Z_{B_4}\) are \(-1/192,-1/96,1/18,1/144\), respectively. These follow from the Euler-product transformation in the radial calculation by factoring the characteristic polynomials above (\((x-1)(x^3-1)\) for \(W\)); a cyclotomic factor \(\Phi_j(x)=\prod_{d|j}(x^d-1)^{m_d}\) contributes \(\sum_{d|j}m_d/(24d)\). The invariant \(C\)-sum has zero exponential rate. Substitution into \eqref{c4:PS-2} proves \eqref{c4:PS-1}, as required.

